% Secciones 4--7 del delta gravitatorio íntegro.
\section{Incidencia hexada--octada, parámetro de Barbero--Immirzi y espectro de área}
\label{sec:l9s102:barbero-area}

La derivación areal recibe como antecedentes la inclusión de incidencia
\(6\to8\), los cardinales TPK \(90/120\), el defecto normalizado
\(-1/12\) y los canales espectrales \(q_\pm\). El funcional
hexada--octada determina el parámetro de Barbero--Immirzi y la
normalización del operador de área. Esta residencia contiene la derivación
HMT y el espectro; la realización Palatini--Cartan--Holst ocupa una residencia
posterior en Mecánica Dimensional.

Sean \(H\subset O\), con \(|H|=6\) y \(|O|=8\), y consérvese el par
marcado de la hexada en la inclusión
\begin{equation}
 \iota_{6,8}:\mathcal F_6=H\times\binom H2
 \lhook\joinrel\longrightarrow
 \mathcal F_8=O\times\binom H2.
 \label{eq:l9s102:barbero-inclusion-incidencia}
\end{equation}
La cuenta interna da
\begin{equation}
 |\mathcal F_6|=6\binom62=90,
 \qquad |\mathcal F_8|=8\binom62=120,
 \qquad
 \frac{90-120}{24\binom62}=-\frac1{12}.
 \label{eq:l9s102:barbero-conteos}
\end{equation}
Una vez publicados \(\pi_{\rm HMT}\), el ángulo \(A\) y el
contraángulo único \(C^*\), los dos canales son
\begin{equation}
 q_\pm=\exp\!\left[-\frac{\pi_{\rm HMT}}{180}(A\pm C^*)\right],
 \qquad 0<q_+<q_-<1,
 \label{eq:l9s102:barbero-qpm}
\end{equation}
y sus series de retorno orientado son
\begin{equation}
 S_{90}(q)=\frac{q^{90}}{1-q^{270}},\qquad
 S_{120}(q)=\frac{q^{120}}{1-q^{360}},\qquad
 \Delta S_m=S_m(q_-)-S_m(q_+).
 \label{eq:l9s102:barbero-series}
\end{equation}

El orden causal de los coeficientes se fija en el registro temporal del TPK,
antes de evaluar el coeficiente de polo. Sea \(C_{12}=\mathbb Z/12\mathbb Z\) la rueda
dodecafásica y sea \(\partial_{\rm ret}\) su única costura orientada de retorno.
En el módulo libre de sucesos de incidencia, la vuelta fundamental es
\begin{equation}
 c_{12}^{+}
 =\sum_{j\in C_{12}}[j,\mathcal F_6]
  -[\partial_{\rm ret},\mathcal F_8].
 \label{eq:l9s102:barbero-cadena-dodecafasica}
\end{equation}
El signo del segundo término es el signo de frontera. La involución
bidireccional invierte la orientación completa,
\(c_{12}^{-}=-c_{12}^{+}\), sin modificar las multiplicidades absolutas de
los dos tipos de incidencia. Defínase el lector de retorno sobre los
generadores por
\begin{equation}
 \mathscr R([j,\mathcal F_6])=S_{90},
 \qquad
 \mathscr R([\partial_{\rm ret},\mathcal F_8])=S_{120}.
 \label{eq:l9s102:barbero-lector-retorno}
\end{equation}
La covariancia de fase hace constante la primera imagen sobre los doce
sectores; la segunda actúa únicamente en la costura ampliada. Por linealidad,
\begin{equation}
 \boxed{\mathscr R(c_{12}^{+})=12S_{90}-S_{120}.}
 \label{eq:l9s102:barbero-peso-generado}
\end{equation}
Así, el peso \(12:-1\) es la imagen de la vuelta fundamental: doce sectores y
una frontera orientada. No procede de imponer previamente \(1/24\), de una
comparación decimal ni del valor convencional de Barbero--Immirzi.

\Needspace{12\baselineskip}
\begin{theorem}[Determinación dodecafásica del funcional de Barbero--Immirzi]
\label{thm:l9s102:barbero-funcional}
La cadena fundamental \eqref{eq:l9s102:barbero-cadena-dodecafasica} determina
el par orientado \((a,b)=(12,1)\). El coeficiente de \((1-q)^{-1}\) de su imagen es
entonces \(1/24\), como salida exacta del funcional ya construido. Por tanto,
después de evaluar los canales conjugados \(q_\pm\), la salida adimensional es
\begin{equation}
 \boxed{
 \gamma_{\rm HMT}=\Gamma_{6\to8}
 =\frac{\pi_{\rm HMT}AC^*}{180}
   +12\Delta S_{90}-\Delta S_{120}.}
 \label{eq:l9s102:barbero-gamma}
\end{equation}
Esta fórmula recibe el par angular ya generado y no utiliza como entrada
el valor convencional del parámetro de Barbero--Immirzi.
\end{theorem}

\begin{proof}
La suma de \eqref{eq:l9s102:barbero-cadena-dodecafasica} contiene exactamente
los doce sectores de \(C_{12}\), mientras la costura de retorno aparece una
sola vez y con orientación negativa. El lector
\eqref{eq:l9s102:barbero-lector-retorno} produce por ello
\eqref{eq:l9s102:barbero-peso-generado}, antes de todo cálculo de polo. Para
cada \(m>0\),
\[
 p_1(S_m)
 =\lim_{q\to1}(1-q)\frac{q^m}{1-q^{3m}}
 =\frac1{3m}.
\]
En consecuencia,
\[
 p_1\bigl(\mathscr R(c_{12}^{+})\bigr)
 =\frac{12}{270}-\frac1{360}
 =\frac{48-3}{1080}=\frac1{24}.
\]
El coeficiente \(p_1\) está referido a \((1-q)^{-1}\).
En la coordenada compleja \(q-1\),
\[
 \operatorname*{Res}_{q=1}S_m(q)=-\frac1{3m},
 \qquad
 \operatorname*{Res}_{q=1}\mathscr R(c_{12}^{+})(q)=-\frac1{24}.
\]
La ecuación diofántica posterior
\[
 \frac a{270}-\frac b{360}=\frac1{24}
 \quad\Longleftrightarrow\quad 4a-3b=45
\]
posee las soluciones enteras positivas
\((a,b)=(12+3t,1+4t)\), \(t\ge0\); la única solución mínima es
\((12,1)\). Esta cuenta certifica aritméticamente el par que la órbita
fundamental ya había generado: cualquier \(t>0\) altera la multiplicidad de
los sectores o de la costura y deja de representar una vuelta del TPK.
Finalmente, sustituir \(q_\pm\) en la imagen orientada y añadir la componente
angular previamente generada prueba \eqref{eq:l9s102:barbero-gamma}.
\end{proof}

La secuencia causal queda fijada sin intercambio de papeles:
\begin{equation}
\begin{aligned}
 (\mathcal F_6\hookrightarrow\mathcal F_8,C_{12},\partial_{\rm ret})
 &\longrightarrow (90,120,-1/12,c_{12}^{+})\\
 &\longrightarrow \mathscr R(c_{12}^{+})=12S_{90}-S_{120},\\
 \mathscr R(c_{12}^{+})
 &\longrightarrow p_1\bigl(\mathscr R(c_{12}^{+})\bigr)=\frac1{24},\\
 \bigl(\mathscr R(c_{12}^{+}),A,C^*,q_\pm\bigr)
 &\longrightarrow \Gamma_{6\to8}=\gamma_{\rm HMT}
 \longrightarrow \widehat{\mathcal A}_{\rm phys}.
\end{aligned}
 \label{eq:l9s102:barbero-cadena-completa}
\end{equation}
El par angular y sus canales proceden de la construcción anterior.
La evaluación del coeficiente de polo constituye una comprobación
del funcional de retorno; la evaluación conjugada emplea las
coordenadas angulares previamente generadas.
La realización posterior en la acción de Holst reconoce la misma sección
adimensional; no interviene en su selección.

Para los dos sectores de borde, sean
\begin{equation}
 N_m=\sum_{e\in\partial A_m}N_e,
 \qquad m\in\{90,120\},
 \qquad \operatorname{Spec}(N_e)=\{0,1,2\}.
 \label{eq:l9s102:n90-n120}
\end{equation}
El operador físico de área se escribe
\begin{equation}
 \boxed{
 \frac{\widehat{\mathcal A}_{\mathrm{phys}}}{\ell_P^2}
 =\gamma_{\mathrm{HMT}}a_0(N_{90}+N_{120}).}
 \label{eq:l9s102:area-operador}
\end{equation}
Sobre \(36\) enlaces qutrit,
\begin{equation}
 \operatorname{Spec}(N_{90}+N_{120})=\{0,1,\ldots,72\},
 \label{eq:l9s102:area-numero-spectrum}
\end{equation}
y la multiplicidad del autovalor \(n\) es el coeficiente de \(z^n\) en
\((1+z+z^2)^{36}\). De aqui resulta
\begin{equation}
 \operatorname{Spec}\!\left(
 \frac{\widehat{\mathcal A}_{\mathrm{phys}}}{\ell_P^2}\right)
 =\{n\gamma_{\mathrm{HMT}}a_0:0\le n\le72\}.
 \label{eq:l9s102:area-spectrum}
\end{equation}

\section{Geometría total y memoria de procedencia}
\label{sec:l9s102:area-hamiltoniano}

El Hamiltoniano modular del borde es
\begin{equation}
 \boxed{
 K_A=-\log\rho_A
 =cI+\Delta_{\mathrm{mod}}(90N_{90}+120N_{120}).}
 \label{eq:l9s102:hamiltoniano-modular}
\end{equation}
Sean
\begin{equation}
 N=N_{90}+N_{120},
 \qquad
 D=90N_{90}+120N_{120}.
 \label{eq:l9s102:n-d}
\end{equation}
La matriz de paso
\(\left(\begin{smallmatrix}1&1\\90&120\end{smallmatrix}\right)\)
posee determinante \(30\), y su inversa reconstruye los sectores:
\begin{equation}
 \boxed{
 N_{90}=4N-\frac{D}{30},
 \qquad
 N_{120}=\frac{D}{30}-3N.}
 \label{eq:l9s102:reconstruccion-sectores}
\end{equation}
El área publica el número total de excitaciones; el Hamiltoniano modular
publica su sector de procedencia. El par conjunto reconstruye la ocupación
completa de la frontera.

\section{Estado bicapa e información orientada}
\label{sec:l9s102:estado-bicapa}

Sea \(R=R^*=R^{-1}\) la involución de hojas y
\(y=C^*_{\mathrm{rad}}\). El estado normalizado es
\begin{equation}
 \boxed{
 \rho_y=\frac{e^{-yR}}{2\cosh y}.}
 \label{eq:l9s102:rho-y}
\end{equation}
Su Hamiltoniano modular y su entropía son
\begin{equation}
 K_y=-\log\rho_y=\log(2\cosh y)I+yR,
 \label{eq:l9s102:ky}
\end{equation}
\begin{equation}
 S(\rho_y)=\log(2\cosh y)-y\tanh y.
 \label{eq:l9s102:entropia-y}
\end{equation}
La inversión \(y\mapsto-y\) deja fija la entropía y produce la distancia
informacional orientada
\begin{equation}
 \boxed{
 D(\rho_y\Vert\rho_{-y})=2y\tanh y.}
 \label{eq:l9s102:entropia-relativa-hojas}
\end{equation}
La componente par mide la distribucion espectral; la componente impar mide
el coste informacional de invertir la orientación.

\section{Cuanto informacional de energía y celda gravitatoria de media acción}
\label{sec:l9s102:horizonte-media-accion}

La periodicidad de \(108\) estados determina
\begin{equation}
 t_P=\frac{54}{\pi}t_0,
 \qquad
 E_P=\frac{\hbar}{t_P}=\frac{h}{108t_0}.
 \label{eq:l9s102:ep-tp}
\end{equation}
La realización térmica del reloj tritico satisface
\begin{equation}
 \boxed{
 E_P=k_B\Theta_{\mathrm{clk}}\log3.}
 \label{eq:l9s102:planck-trit}
\end{equation}
Para el horizonte de radio \(R\),
\begin{equation}
 E_\Lambda(R)=\frac{c^4R}{2G},
 \qquad
 E_\Lambda(R)=T_H(R)S_H(R).
 \label{eq:l9s102:energia-horizonte}
\end{equation}
En la celda de Planck,
\begin{equation}
 \boxed{
 E_\Lambda
 =T_HS_H
 =\frac{E_P}{2}
 =\frac{k_B\Theta_{\mathrm{clk}}\log3}{2}.}
 \label{eq:l9s102:confluencia-horizonte}
\end{equation}
Por tanto,
\begin{equation}
 \boxed{
 E_\Lambda t_P=\frac{\hbar}{2},
 \qquad
 E_\Lambda T_{108}=\frac h2.}
 \label{eq:l9s102:media-accion}
\end{equation}
La energía de vacío, la entropía areal, la temperatura cronológica y la
acción comparten una misma celda de confluencia con dominios tipados.

Cuando \(S_H/k_B=\pi\), la multiplicidad efectiva satisface
\begin{equation}
 \Omega_{\mathrm{eff}}=e^{S_H/k_B}=e^\pi.
 \label{eq:l9s102:omega-e-pi}
\end{equation}
El operador hiperbólico de Euler posee
\begin{equation}
 \operatorname{Spec}(e^{\pi H})=\{e^\pi,e^{-\pi}\}.
 \label{eq:l9s102:e-pi-spectrum}
\end{equation}
La igualdad de los escalares establece la confluencia espectral; el
entrelazador entre multiplicidad microscópica de horizonte y dinamica de
Perron pertenece a su residencia física especifica.
